% --- Recap / bridge from earlier lectures ---
\begin{frame}{Recap: From Pretraining Methods to Foundation Models}
  \begin{itemize}\itemsep0.3em
    \item \textbf{Covered earlier:} pretext tasks, contrastive learning (SimCLR, MoCo), and self-distillation (BYOL, DINO, DINOv2) learn features from unlabeled data.
    \item \textbf{Also covered earlier:} latent prediction (I-JEPA, V-JEPA 2) trains encoders that serve many tasks from frozen features; V-JEPA 2 already behaves like a video foundation model.
    \item Those lectures answered \emph{how to pretrain}. Today we look at the pretrained models the field actually builds on: what they are, what data they use, and how they are adapted.
    \item Working definition (Bommasani et al., 2021): a \textbf{foundation model} is trained on broad data at scale and can be \textbf{adapted} to a wide range of downstream tasks.
  \end{itemize}
\end{frame}
